\chapter*{Checklista wymagań z szablonu}

\section*{Strona tytułowa i układ dokumentu}
\begin{itemize}
    \item Praca ma być złożona jako praca dyplomowa magisterska w języku polskim.
    \item Strona tytułowa powinna zawierać autorów z numerami albumów, tytuł pracy, promotora oraz rok złożenia pracy.
    \item Tytuł pracy należy sprawdzić wizualnie: powinien mieć maksymalnie dwa wiersze i nie powinien być łamany w nielogicznym miejscu.
    \item Dokument używa formatu A4, układu jednostronnego, marginesów zdefiniowanych w klasie oraz interlinii półtorej.
    \item Praca składa się z rozdziałów i podrozdziałów; nie należy schodzić głębiej niż do poziomu \texttt{subsubsection}.
\end{itemize}

\section*{Wstęp}
\begin{itemize}
    \item Wstęp powinien krótko uzasadniać podjęcie tematu.
    \item We wstępie należy podać cel pracy.
    \item We wstępie należy określić zakres pracy: przedmiotowy, podmiotowy i czasowy, jeżeli ma zastosowanie.
    \item Można wskazać hipotezy, które autorzy zamierzają sprawdzić lub udowodnić.
    \item Należy krótko scharakteryzować źródła, zwłaszcza literaturowe.
    \item Należy opisać układ pracy, czyli zawartość poszczególnych rozdziałów.
    \item Można dodać uwagi dotyczące realizacji tematu, np. trudności, sprzętu lub współpracy z podmiotami zewnętrznymi.
    \item Wstęp musi zawierać akapit rozpoczynający się od słów: „Celem pracy jest ...”.
    \item Wstęp musi zawierać akapit rozpoczynający się od słów: „Struktura pracy jest następująca. ...”.
    \item W przypadku dwuosobowej pracy magisterskiej po tych akapitach należy opisać wkład poszczególnych członków zespołu.
\end{itemize}

\section*{Podstawy teoretyczne i literatura}
\begin{itemize}
    \item Rozdział teoretyczny powinien stanowić przegląd literatury naświetlający aktualny stan wiedzy.
    \item W tekście pracy muszą wystąpić odwołania do wszystkich pozycji z wykazu literatury.
    \item Odnośników do literatury nie należy umieszczać w stopkach stron.
    \item Należy wskazywać źródła pochodzenia informacji, rysunków, tabel i fragmentów kodu.
    \item Dla źródeł internetowych należy podać adresy stron.
\end{itemize}

\section*{Praca własna, implementacja i wyniki}
\begin{itemize}
    \item Praca musi zawierać elementy pracy własnej autorów adekwatne do wiedzy praktycznej uzyskanej w czasie studiów.
    \item Za pracę własną można uznać m.in. aplikację, fragment aplikacji, algorytm, projekt systemu, analizę lub ocenę technologii.
    \item Należy udokumentować założenia, sposób realizacji zadań, ocenę wykonania i napotkane problemy.
    \item Dla pracy projektowo-implementacyjnej należy przygotować dokumentację techniczną i użytkową systemu.
    \item Nie należy zamieszczać w pracy całego kodu źródłowego.
    \item Kod, oprogramowanie i wyniki eksperymentów powinny zostać umieszczone w dodatku lub innym załączniku.
\end{itemize}

\section*{Styl tekstu}
\begin{itemize}
    \item Należy stosować formę bezosobową, np. „w pracy rozważono” albo „w ramach pracy zaprojektowano”.
    \item Należy unikać długich zdań.
    \item Nie należy używać zwrotów potocznych.
    \item Należy używać precyzyjnej terminologii informatycznej.
    \item Nie wolno pisać pracy metodą kopiuj-wklej; zagadnienia trzeba opisywać własnymi słowami.
    \item Przy informacjach pochodzących z zewnątrz zawsze należy wskazywać źródło.
\end{itemize}

\section*{Zakończenie}
\begin{itemize}
    \item Zakończenie powinno odnosić się do zadań wskazanych we wstępie.
    \item Należy odnieść się szczególnie do celu i zakresu pracy.
    \item Należy porównać założenia z faktycznymi wynikami pracy.
    \item Należy jasno określić stopień realizacji celów oraz wyniki samodzielnej pracy autorów.
    \item Dodatki, aneksy i załączniki są integralną częścią pracy.
\end{itemize}

\section*{Zasady składu w LaTeX}
\begin{itemize}
    \item Akapity rozdziela się co najmniej jedną pustą linią.
    \item Wyróżnienia należy wykonywać poleceniem \texttt{\textbackslash emph}; pogrubienie nie jest zalecane jako zwykłe wyróżnienie.
    \item Po jednoliterowych spójnikach warto stosować twardą spację, czyli znak tyldy.
    \item Należy używać poprawnych polskich cudzysłowów i konsekwentnego zapisu pauz oraz półpauz.
    \item Rysunki należy umieszczać w środowisku \texttt{figure}, z podpisem pod rysunkiem.
    \item Rysunki powinny być umieszczane u góry strony; nie zaleca się modyfikatora \texttt{[h]}.
    \item Tabele należy umieszczać w środowisku \texttt{table}, z podpisem nad tabelą.
    \item Podpisy rysunków i tabel powinny kończyć się kropką.
    \item Można używać makr \texttt{\textbackslash definicja}, \texttt{\textbackslash akronim} i \texttt{\textbackslash english}.
    \item Nazwy plików, katalogów, ścieżek, zmiennych, klas i metod należy formatować poleceniem \texttt{\textbackslash texttt}.
    \item Przypisów nie należy używać zbyt często; warto rozważyć włączenie treści przypisu do tekstu głównego.
    \item Po zmianach należy sprawdzić ostrzeżenia \texttt{overfull} i \texttt{underfull boxes}.
    \item Po ostatniej redakcji należy sprawdzić wiszące spójniki i wstawić twarde spacje tam, gdzie są potrzebne.
\end{itemize}
